\IfFileExists{stacks-project.cls}{%
\documentclass{stacks-project}
}{%
\documentclass{amsart}
}

% For dealing with references we use the comment environment
\usepackage{verbatim}
\newenvironment{reference}{\comment}{\endcomment}
%\newenvironment{reference}{}{}
\newenvironment{slogan}{\comment}{\endcomment}
\newenvironment{history}{\comment}{\endcomment}

% For commutative diagrams we use Xy-pic
\usepackage[all]{xy}

% We use 2cell for 2-commutative diagrams.
\xyoption{2cell}
\UseAllTwocells

% We use multicol for the list of chapters between chapters
\usepackage{multicol}

% This is generally recommended for better output
\usepackage{lmodern}
\usepackage[T1]{fontenc}

% For cross-file-references
\usepackage{xr-hyper}

% Package for hypertext links:
\usepackage{hyperref}

% For any local file, say "hello.tex" you want to link to please
% use \externaldocument[hello-]{hello}
\externaldocument[introduction-]{introduction}
\externaldocument[conventions-]{conventions}
\externaldocument[sets-]{sets}
\externaldocument[categories-]{categories}
\externaldocument[topology-]{topology}
\externaldocument[sheaves-]{sheaves}
\externaldocument[sites-]{sites}
\externaldocument[stacks-]{stacks}
\externaldocument[fields-]{fields}
\externaldocument[algebra-]{algebra}
\externaldocument[brauer-]{brauer}
\externaldocument[homology-]{homology}
\externaldocument[derived-]{derived}
\externaldocument[simplicial-]{simplicial}
\externaldocument[more-algebra-]{more-algebra}
\externaldocument[smoothing-]{smoothing}
\externaldocument[modules-]{modules}
\externaldocument[sites-modules-]{sites-modules}
\externaldocument[injectives-]{injectives}
\externaldocument[cohomology-]{cohomology}
\externaldocument[sites-cohomology-]{sites-cohomology}
\externaldocument[dga-]{dga}
\externaldocument[dpa-]{dpa}
\externaldocument[sdga-]{sdga}
\externaldocument[hypercovering-]{hypercovering}
\externaldocument[schemes-]{schemes}
\externaldocument[constructions-]{constructions}
\externaldocument[properties-]{properties}
\externaldocument[morphisms-]{morphisms}
\externaldocument[coherent-]{coherent}
\externaldocument[divisors-]{divisors}
\externaldocument[limits-]{limits}
\externaldocument[varieties-]{varieties}
\externaldocument[topologies-]{topologies}
\externaldocument[descent-]{descent}
\externaldocument[perfect-]{perfect}
\externaldocument[more-morphisms-]{more-morphisms}
\externaldocument[flat-]{flat}
\externaldocument[groupoids-]{groupoids}
\externaldocument[more-groupoids-]{more-groupoids}
\externaldocument[etale-]{etale}
\externaldocument[chow-]{chow}
\externaldocument[intersection-]{intersection}
\externaldocument[pic-]{pic}
\externaldocument[weil-]{weil}
\externaldocument[adequate-]{adequate}
\externaldocument[dualizing-]{dualizing}
\externaldocument[duality-]{duality}
\externaldocument[discriminant-]{discriminant}
\externaldocument[derham-]{derham}
\externaldocument[local-cohomology-]{local-cohomology}
\externaldocument[algebraization-]{algebraization}
\externaldocument[curves-]{curves}
\externaldocument[resolve-]{resolve}
\externaldocument[models-]{models}
\externaldocument[functors-]{functors}
\externaldocument[equiv-]{equiv}
\externaldocument[pione-]{pione}
\externaldocument[etale-cohomology-]{etale-cohomology}
\externaldocument[proetale-]{proetale}
\externaldocument[relative-cycles-]{relative-cycles}
\externaldocument[more-etale-]{more-etale}
\externaldocument[trace-]{trace}
\externaldocument[crystalline-]{crystalline}
\externaldocument[spaces-]{spaces}
\externaldocument[spaces-properties-]{spaces-properties}
\externaldocument[spaces-morphisms-]{spaces-morphisms}
\externaldocument[decent-spaces-]{decent-spaces}
\externaldocument[spaces-cohomology-]{spaces-cohomology}
\externaldocument[spaces-limits-]{spaces-limits}
\externaldocument[spaces-divisors-]{spaces-divisors}
\externaldocument[spaces-over-fields-]{spaces-over-fields}
\externaldocument[spaces-topologies-]{spaces-topologies}
\externaldocument[spaces-descent-]{spaces-descent}
\externaldocument[spaces-perfect-]{spaces-perfect}
\externaldocument[spaces-more-morphisms-]{spaces-more-morphisms}
\externaldocument[spaces-flat-]{spaces-flat}
\externaldocument[spaces-groupoids-]{spaces-groupoids}
\externaldocument[spaces-more-groupoids-]{spaces-more-groupoids}
\externaldocument[bootstrap-]{bootstrap}
\externaldocument[spaces-pushouts-]{spaces-pushouts}
\externaldocument[spaces-chow-]{spaces-chow}
\externaldocument[groupoids-quotients-]{groupoids-quotients}
\externaldocument[spaces-more-cohomology-]{spaces-more-cohomology}
\externaldocument[spaces-simplicial-]{spaces-simplicial}
\externaldocument[spaces-duality-]{spaces-duality}
\externaldocument[formal-spaces-]{formal-spaces}
\externaldocument[restricted-]{restricted}
\externaldocument[spaces-resolve-]{spaces-resolve}
\externaldocument[formal-defos-]{formal-defos}
\externaldocument[defos-]{defos}
\externaldocument[cotangent-]{cotangent}
\externaldocument[examples-defos-]{examples-defos}
\externaldocument[algebraic-]{algebraic}
\externaldocument[examples-stacks-]{examples-stacks}
\externaldocument[stacks-sheaves-]{stacks-sheaves}
\externaldocument[criteria-]{criteria}
\externaldocument[artin-]{artin}
\externaldocument[quot-]{quot}
\externaldocument[stacks-properties-]{stacks-properties}
\externaldocument[stacks-morphisms-]{stacks-morphisms}
\externaldocument[stacks-limits-]{stacks-limits}
\externaldocument[stacks-cohomology-]{stacks-cohomology}
\externaldocument[stacks-perfect-]{stacks-perfect}
\externaldocument[stacks-introduction-]{stacks-introduction}
\externaldocument[stacks-more-morphisms-]{stacks-more-morphisms}
\externaldocument[stacks-geometry-]{stacks-geometry}
\externaldocument[moduli-]{moduli}
\externaldocument[moduli-curves-]{moduli-curves}
\externaldocument[examples-]{examples}
\externaldocument[exercises-]{exercises}
\externaldocument[guide-]{guide}
\externaldocument[desirables-]{desirables}
\externaldocument[coding-]{coding}
\externaldocument[obsolete-]{obsolete}
\externaldocument[fdl-]{fdl}
\externaldocument[index-]{index}

% Theorem environments.
%
\theoremstyle{plain}
\newtheorem{theorem}[subsection]{Theorem}
\newtheorem{proposition}[subsection]{Proposition}
\newtheorem{lemma}[subsection]{Lemma}

\theoremstyle{definition}
\newtheorem{definition}[subsection]{Definition}
\newtheorem{example}[subsection]{Example}
\newtheorem{exercise}[subsection]{Exercise}
\newtheorem{situation}[subsection]{Situation}

\theoremstyle{remark}
\newtheorem{remark}[subsection]{Remark}
\newtheorem{remarks}[subsection]{Remarks}

\numberwithin{equation}{subsection}

% Macros
%
\def\lim{\mathop{\mathrm{lim}}\nolimits}
\def\colim{\mathop{\mathrm{colim}}\nolimits}
\def\Spec{\mathop{\mathrm{Spec}}}
\def\Hom{\mathop{\mathrm{Hom}}\nolimits}
\def\Ext{\mathop{\mathrm{Ext}}\nolimits}
\def\SheafHom{\mathop{\mathcal{H}\!\mathit{om}}\nolimits}
\def\SheafExt{\mathop{\mathcal{E}\!\mathit{xt}}\nolimits}
\def\Sch{\mathit{Sch}}
\def\Mor{\mathop{\mathrm{Mor}}\nolimits}
\def\Ob{\mathop{\mathrm{Ob}}\nolimits}
\def\Sh{\mathop{\mathit{Sh}}\nolimits}
\def\NL{\mathop{N\!L}\nolimits}
\def\CH{\mathop{\mathrm{CH}}\nolimits}
\def\proetale{{pro\text{-}\acute{e}tale}}
\def\etale{{\acute{e}tale}}
\def\QCoh{\mathit{QCoh}}
\def\Ker{\mathop{\mathrm{Ker}}}
\def\Im{\mathop{\mathrm{Im}}}
\def\Coker{\mathop{\mathrm{Coker}}}
\def\Coim{\mathop{\mathrm{Coim}}}

% Boxtimes
%
\DeclareMathSymbol{\boxtimes}{\mathbin}{AMSa}{"02}

%
% Macros for moduli stacks/spaces
%
\def\QCohstack{\mathcal{QC}\!\mathit{oh}}
\def\Cohstack{\mathcal{C}\!\mathit{oh}}
\def\Spacesstack{\mathcal{S}\!\mathit{paces}}
\def\Quotfunctor{\mathrm{Quot}}
\def\Hilbfunctor{\mathrm{Hilb}}
\def\Curvesstack{\mathcal{C}\!\mathit{urves}}
\def\Polarizedstack{\mathcal{P}\!\mathit{olarized}}
\def\Complexesstack{\mathcal{C}\!\mathit{omplexes}}
% \Pic is the operator that assigns to X its picard group, usage \Pic(X)
% \Picardstack_{X/B} denotes the Picard stack of X over B
% \Picardfunctor_{X/B} denotes the Picard functor of X over B
\def\Pic{\mathop{\mathrm{Pic}}\nolimits}
\def\Picardstack{\mathcal{P}\!\mathit{ic}}
\def\Picardfunctor{\mathrm{Pic}}
\def\Deformationcategory{\mathcal{D}\!\mathit{ef}}
\usepackage{cleveref}
\crefname{originallemma}{original lemma}{original lemmas}
\Crefname{originallemma}{Original lemma}{Original lemmas}
\makeatletter
\newlabel{etale-cohomology-lemma-jshriek-open}{{0F70}{source}{}{sourcepointer-1}{}}
\newlabel{etale-cohomology-lemma-jshriek-open@cref}{{[originallemma][0][]0F70}{[1][0][]source}{}{}{}}
\makeatother
\crefname{originallemma}{original lemma}{original lemmas}
\Crefname{originallemma}{Original lemma}{Original lemmas}
\makeatletter
\newlabel{etale-cohomology-lemma-proper-pushforward-stalk}{{0A3T}{source}{}{sourcepointer-2}{}}
\newlabel{etale-cohomology-lemma-proper-pushforward-stalk@cref}{{[originallemma][0][]0A3T}{[1][0][]source}{}{}{}}
\makeatother
\crefname{originaltheorem}{original theorem}{original theorems}
\Crefname{originaltheorem}{Original theorem}{Original theorems}
\makeatletter
\newlabel{flat-theorem-nagata}{{0F41}{source}{}{sourcepointer-3}{}}
\newlabel{flat-theorem-nagata@cref}{{[originaltheorem][0][]0F41}{[1][0][]source}{}{}{}}
\makeatother
\crefname{originallemma}{original lemma}{original lemmas}
\Crefname{originallemma}{Original lemma}{Original lemmas}
\makeatletter
\newlabel{lemma-compactify-compact-support}{{0F75}{source}{}{sourcepointer-4}{}}
\newlabel{lemma-compactify-compact-support@cref}{{[originallemma][0][]0F75}{[1][0][]source}{}{}{}}
\makeatother
\crefname{originallemma}{original lemma}{original lemmas}
\Crefname{originallemma}{Original lemma}{Original lemmas}
\makeatletter
\newlabel{lemma-compactify-f-shriek-separated}{{0F52}{source}{}{sourcepointer-5}{}}
\newlabel{lemma-compactify-f-shriek-separated@cref}{{[originallemma][0][]0F52}{[1][0][]source}{}{}{}}
\makeatother
\crefname{originallemma}{original lemma}{original lemmas}
\Crefname{originallemma}{Original lemma}{Original lemmas}
\makeatletter
\newlabel{lemma-proper-compact-support}{{0F73}{source}{}{sourcepointer-6}{}}
\newlabel{lemma-proper-compact-support@cref}{{[originallemma][0][]0F73}{[1][0][]source}{}{}{}}
\makeatother
\crefname{originallemma}{original lemma}{original lemmas}
\Crefname{originallemma}{Original lemma}{Original lemmas}
\makeatletter
\newlabel{lemma-proper-f-shriek}{{0F51}{source}{}{sourcepointer-7}{}}
\newlabel{lemma-proper-f-shriek@cref}{{[originallemma][0][]0F51}{[1][0][]source}{}{}{}}
\makeatother
\crefname{originalremark}{original remark}{original remarks}
\Crefname{originalremark}{Original remark}{Original remarks}
\makeatletter
\newlabel{remark-covariance-compact-support}{{0F74}{source}{}{sourcepointer-8}{}}
\newlabel{remark-covariance-compact-support@cref}{{[originalremark][0][]0F74}{[1][0][]source}{}{}{}}
\makeatother
\crefname{originalremark}{original remark}{original remarks}
\Crefname{originalremark}{Original remark}{Original remarks}
\makeatletter
\newlabel{remark-covariance-f-shriek-separated}{{0F53}{source}{}{sourcepointer-9}{}}
\newlabel{remark-covariance-f-shriek-separated@cref}{{[originalremark][0][]0F53}{[1][0][]source}{}{}{}}
\makeatother
\title{Stalks of lower shriek identified with compact-support sections on the geometric fibre}
\author{The Stacks Project Authors}
\date{stacks-project:more-etale.tex:local-snapshot-20261012; GNU FDL 1.2 or later}
\begin{document}
\maketitle
\setcounter{subsection}{1}
\noindent\textit{Editorial note.} All delivered mathematical text is the author's own, copied from the pinned Stacks Project chapter file more-etale.tex; no statement or proof step was rewritten, reordered or removed. The counted span is the author environment and its own proof environment. References inside the span are handed to the assembler as navigation stubs that carry the canonical tag of the cited tag only; whether each of them is load-bearing is recorded on the stub itself.\par\medskip
\noindent\textit{Allowed standard prerequisites (editorial).}\begin{itemize}
\item Nagata compactification (More on Flatness, Theorem flat-theorem-nagata, canonical tag 0F41). cited at proof lines 737-739. Provides the compactification used to reduce the quasi-compact affine case to the proper case.
\item Proper base change for stalks of pushforwards (Etale Cohomology, Lemma etale-cohomology-lemma-proper-pushforward-stalk, canonical tag 0A3T). cited at proof lines 757-759. Identifies the stalk of Rf\_* for the proper morphism bar-X -> Y in the final reduction.
\end{itemize}\medskip
\ifdefined\hypertarget\hypertarget{sourcepointer-1}{}\fi
\noindent\textit{Source reference.} Original lemma 0F70: sources/stacks-chow/tags.txt, source lines 18990--18990. Cited at proof line 741: F = G|\_X for G = j\_!F in the compactification reduction.\par\medskip
\ifdefined\hypertarget\hypertarget{sourcepointer-2}{}\fi
\noindent\textit{Source reference.} Original lemma 0A3T: sources/stacks-chow/tags.txt, source lines 12751--12751. Cited at proof lines 757-759: identifies the stalk of R(bar-f)\_* G in the proper case.\par\medskip
\ifdefined\hypertarget\hypertarget{sourcepointer-3}{}\fi
\noindent\textit{Source reference.} Original theorem 0F41: sources/stacks-chow/tags.txt, source lines 18886--18886. Cited at proof lines 737-739: existence of the compactification j: X -> bar-X over Y.\par\medskip
\ifdefined\hypertarget\hypertarget{sourcepointer-4}{}\fi
\noindent\textit{Source reference.} Original lemma 0F75: sources/stacks-chow/tags.txt, source lines 18995--18995. Cited at proof lines 751-752: the right vertical map of the reduction square is an isomorphism.\par\medskip
\ifdefined\hypertarget\hypertarget{sourcepointer-5}{}\fi
\noindent\textit{Source reference.} Original lemma 0F52: sources/stacks-chow/tags.txt, source lines 18922--18922. Cited at proof lines 751-752: the left vertical map of the reduction square is an isomorphism.\par\medskip
\ifdefined\hypertarget\hypertarget{sourcepointer-6}{}\fi
\noindent\textit{Source reference.} Original lemma 0F73: sources/stacks-chow/tags.txt, source lines 18993--18993. Cited at proof line 756: compact-support sections equal sections over the fibre in the proper case.\par\medskip
\ifdefined\hypertarget\hypertarget{sourcepointer-7}{}\fi
\noindent\textit{Source reference.} Original lemma 0F51: sources/stacks-chow/tags.txt, source lines 18921--18921. Cited at proof line 756: lower shriek equals pushforward for the proper morphism bar-X -> Y.\par\medskip
\ifdefined\hypertarget\hypertarget{sourcepointer-8}{}\fi
\noindent\textit{Source reference.} Original remark 0F74: sources/stacks-chow/tags.txt, source lines 18994--18994. Cited at proof lines 688-689: the right vertical arrow of the compatibility square.\par\medskip
\ifdefined\hypertarget\hypertarget{sourcepointer-9}{}\fi
\noindent\textit{Source reference.} Original remark 0F53: sources/stacks-chow/tags.txt, source lines 18923--18923. Cited at proof lines 688-689: the left vertical arrow of the compatibility square.\par\medskip
\begin{lemma}
\label{lemma-stalk-f-shriek-separated}
Let $f : X \to Y$ be a morphism of schemes which is separated and
locally of finite type. Let $\mathcal{F}$ be an abelian sheaf on
$X_\etale$. Then there is a canonical isomorphism
$$
(f_!\mathcal{F})_{\overline{y}}
\longrightarrow
H^0_c(X_{\overline{y}}, \mathcal{F}|_{X_{\overline{y}}})
$$
for any geometric point $\overline{y} : \Spec(k) \to Y$.
\end{lemma}

\begin{proof}
Recall that $(f_*\mathcal{F})_{\overline{y}} = \colim f_*\mathcal{F}(V)$
where the colimit is over the \'etale neighbourhoods $(V, \overline{v})$
of $\overline{y}$. If $s \in f_*\mathcal{F}(V) = \mathcal{F}(X_V)$,
then we can pullback $s$ to a section of $\mathcal{F}$ over
$(X_V)_{\overline{v}} = X_{\overline{y}}$. Thus we obtain a canonical map
$$
c_{\overline{y}} :
(f_*\mathcal{F})_{\overline{y}}
\longrightarrow 
H^0(X_{\overline{y}}, \mathcal{F}|_{X_{\overline{y}}})
$$
We claim that this map induces a bijection between the subgroups
$(f_!\mathcal{F})_{\overline{y}}$ and
$H^0_c(X_{\overline{y}}, \mathcal{F}|_{X_{\overline{y}}})$.
The claim implies the lemma, but is a little bit more precise
in that it describes the identification of the lemma as given
by pullbacks of sections of $\mathcal{F}$ to the geometric fibre of $f$.

\medskip\noindent
Observe that any element
$s \in (f_!\mathcal{F})_{\overline{y}} \subset (f_*\mathcal{F})_{\overline{y}}$
is mapped by $c_{\overline{y}}$ to an element of
$H^0_c(X_{\overline{y}}, \mathcal{F}|_{X_{\overline{y}}}) \subset
H^0(X_{\overline{y}}, \mathcal{F}|_{X_{\overline{y}}})$.
This is true because taking the support of a section
commutes with pullback and because properness is preserved by
base change. This at least produces the map in the statement of the lemma.
To prove that it is an isomorphism we may work Zariski
locally on $Y$ and hence we may and do assume $Y$ is affine.

\medskip\noindent
An observation that we will use below
is that given an open subscheme $X' \subset X$
and if $f' = f|_{X'}$, then we obtain a commutative diagram
$$
\xymatrix{
(f'_!(\mathcal{F}|_{X'}))_{\overline{y}} \ar[r] \ar[d] &
H^0_c(X'_{\overline{y}}, \mathcal{F}|_{X'_{\overline{y}}}) \ar[d] \\
(f_!\mathcal{F})_{\overline{y}} \ar[r] &
H^0_c(X_{\overline{y}}, \mathcal{F}|_{X_{\overline{y}}})
}
$$
where the horizontal arrows are the maps constructed above and
the vertical arrows are given in
Remarks \ref{remark-covariance-f-shriek-separated} and
\ref{remark-covariance-compact-support}.
The reason is that given an \'etale neighbourhood $(V, \overline{v})$
of $\overline{y}$ and a section $s \in f_*\mathcal{F}(V) = \mathcal{F}(X_V)$
whose support $Z$ happens to be contained in $X'_V$ and is proper over $V$,
so that $s$ gives rise to an element of both
$(f'_!(\mathcal{F}|_{X'}))_{\overline{y}}$ and
$(f_!\mathcal{F})_{\overline{y}}$ which correspond via
the vertical arrow of the diagram, then these elements are mapped via the
horizontal arrows to the pullback $s|_{X_{\overline{y}}}$ of $s$ to
$X_{\overline{y}}$ whose support $Z_{\overline{y}}$ is contained in
$X'_{\overline{y}}$ and hence this restriction gives rise to
a compatible pair of elements of
$H^0_c(X'_{\overline{y}}, \mathcal{F}|_{X'_{\overline{y}}})$
and
$H^0_c(X_{\overline{y}}, \mathcal{F}|_{X_{\overline{y}}})$.

\medskip\noindent
Suppose $s \in (f_!\mathcal{F})_{\overline{y}}$ maps to zero in
$H^0_c(X_{\overline{y}}, \mathcal{F}|_{X_{\overline{y}}})$.
Say $s$ corresponds to $s \in f_*\mathcal{F}(V) = \mathcal{F}(X_V)$
with support $Z$ proper over $V$. We may assume that $V$ is affine
and hence $Z$ is quasi-compact. Then we may choose a quasi-compact open
$X' \subset X$ containing the image of $Z$. Then $Z$ is contained in
$X'_V$ and hence $s$ is the image of an element
$s' \in f'_!(\mathcal{F}|_{X'})(V)$ where $f' = f|_{X'}$ as in
the previous paragraph. Then $s'$ maps to zero in
$H^0_c(X'_{\overline{y}}, \mathcal{F}|_{X'_{\overline{y}}})$.
Hence in order to prove injectivity, we may replace $X$ by
$X'$, i.e., we may assume $X$ is quasi-compact. We will prove
this case below.

\medskip\noindent
Suppose that
$t \in H^0_c(X_{\overline{y}}, \mathcal{F}|_{X_{\overline{y}}})$.
Then the support of $t$ is contained in a quasi-compact
open subscheme $W \subset X_{\overline{y}}$.
Hence we can find a quasi-compact open subscheme
$X' \subset X$ such that $X'_{\overline{y}}$ contains $W$.
Then it is clear that $t$ is contained in the image
of the injective map
$H^0_c(X'_{\overline{y}}, \mathcal{F}|_{X'_{\overline{y}}}) \to
H^0_c(X_{\overline{y}}, \mathcal{F}|_{X_{\overline{y}}})$.
Hence in order to show surjectivity, we may replace $X$
by $X'$, i.e., we may assume $X$ is quasi-compact.
We will prove this case below.

\medskip\noindent
In this last paragraph of the proof we prove the lemma in case
$X$ is quasi-compact and $Y$ is affine. By More on Flatness, Theorem
\ref{flat-theorem-nagata} there exists a compactification
$j : X \to \overline{X}$ over $Y$. Set $\mathcal{G} = j_!\mathcal{F}$
so that $\mathcal{F} = \mathcal{G}|_X$ by
\'Etale Cohomology, Lemma \ref{etale-cohomology-lemma-jshriek-open}.
By the discussion above we get a commutative diagram
$$
\xymatrix{
(f_!\mathcal{F})_{\overline{y}} \ar[r] \ar[d] &
H^0_c(X_{\overline{y}}, \mathcal{F}|_{X_{\overline{y}}}) \ar[d] \\
(\overline{f}_!\mathcal{G})_{\overline{y}} \ar[r] &
H^0_c(\overline{X}_{\overline{y}}, \mathcal{G}|_{\overline{X}_{\overline{y}}})
}
$$
By Lemmas \ref{lemma-compactify-f-shriek-separated} and
\ref{lemma-compactify-compact-support} the vertical maps
are isomorphisms. This reduces us to the case of the proper
morphism $\overline{X} \to Y$. For a proper morphism our map
is an isomorphism by
Lemmas \ref{lemma-proper-f-shriek} and \ref{lemma-proper-compact-support}
and proper base change for pushforwards, see
\'Etale Cohomology, Lemma
\ref{etale-cohomology-lemma-proper-pushforward-stalk}.
\end{proof}

\end{document}
